\documentclass[10pt,a4paper]{amsart}
\usepackage[margin=19mm]{geometry}
\usepackage{amsmath,amssymb,mathtools}
\usepackage[scr=rsfs,cal=euler]{mathalfa}
\usepackage{xfrac}
\usepackage[unicode,hidelinks,bookmarks=false]{hyperref}
\newcommand{\ie}{i.e.\ }
\newcommand{\cf}{cf.\ }
\newcommand{\resp}{respectively\ }
\newcommand{\Subsection}{\subsection}
\providecommand{\nobreakdash}{\nobreak\hbox{-}}
\setlength{\emergencystretch}{2em}
\multlinegap0pt
\usepackage{amssymb}
\newtheorem{corollary}{Corollary}
\newtheorem{proposition}{Proposition}
\usepackage{cleveref}
\crefname{corollary}{Corollary}{Corollarys}
\crefname{proposition}{Proposition}{Propositions}
\newcommand{\C}{\mathbb{C}}
\newcommand{\End}{\text{End}}
\newcommand{\Hom}{\text{Hom}}
\newcommand{\Q}{\mathbb{Q}}
\newcommand{\R}{\mathbb{R}}
\title{The monodromy of compact Lagrangian fibrations}
\author{Edward Varvak}
\date{Source: arXiv:2504.21740v2 (CC BY 4.0)}
\begin{document}
\maketitle

\noindent\emph{Source and licence.} The monodromy of compact Lagrangian fibrations. Version arXiv:2504.21740v2 (CC BY 4.0), distributed by arXiv under the Creative Commons Attribution 4.0 International licence, http://creativecommons.org/licenses/by/4.0/, as recorded in the arXiv licence metadata for this exact version and shown on the abstract page https://arxiv.org/abs/2504.21740v2. Full licence notice: \texttt{licenses/arxiv-2504.21740-CC-BY-4.0.txt}.\par\smallskip

\noindent\textit{Editorial scope.} Counted unit: the article's proposition isotrivial-rep-classification (source0.tex lines 947--960). The article's complete original proof of it is delivered with it (source0.tex lines 962--1032, one \texttt{proof} environment of 71 lines). No mathematics is written, repaired or re-derived: the statement and the proof are the article's own lines, in the article's own notation, and no retained line is substituted or re-flowed.  Retained uncounted as the source-local context the proof needs: lines 419--424.  Nothing was omitted from any retained span, so the package carries no omission record.  No retained line contains a citation, so the wrapper carries no bibliography and no bibliography entry.  No retained line contains a figure environment, an image-inclusion command or drawing code, so no image is delivered and the package declares no asset dependency.  Editorial formatting only, in addition: an AMS \texttt{amsart} preamble that restates the article's own declarations and macros used by the retained lines, namely the environments \texttt{corollary}, \texttt{proposition} on separate counters, together with the standard packages those lines need.  The retained environments print in this document as Corollary 1, Proposition 1 in order of appearance, whereas the article prints the same items with the numbering of its own full document.  The complete native source of the article as distributed in the arXiv e-print for this exact version is bundled unchanged as \texttt{sources/177359/source0.tex}.
\begin{corollary}\label{proof-of-rep-classification}
    Let \(U\) be an irreducible complex variation of Hodge structures with 
    \(U \simeq \overline{U}\). Then \(h^0(\mathrm{Sym}^2 U) = \C\varphi \neq 0\) 
    if and only if \(U\) is of real type, and \(h^0(\bigwedge^2 U) = \C \varphi \neq 0\)
    if and only if \(U\) is of quaternionic type.
\end{corollary}

\begin{proposition} \label{isotrivial-rep-classification}
    Let \(E\) be the elliptic curve from above, \(H = H^1(E, \Q)\).
    \begin{itemize}
    \item[(a)] \(V_\C\) is of real type if and only if we have an isomorphism
    of \(\Q\)-Hodge representations \(V_\Q = H \otimes W\) for \(W\) a
    \(\Q\)-local system irreducible over \(\C\).
    \item[(b)] Let \(K = \mathrm{End}_{\Q HS}(H)\) be the CM field of \(E\).
    \(V_\C\) is of complex type precisely when \(K \neq \Q\) and there is an 
    isomorphism of \(K\)-Hodge representations  \(V_K \simeq V_K^{1,0} \oplus V_K^{0,1}\).
    Here \(V_K^{1,0}\) and \(V_K^{0,1}\) are non-isomorphic \(K\)-local systems which
    are irreducible over \(\C\).
    \item[(c)] \(V_\C\) is not quaternionic.
    \end{itemize}
\end{proposition}

\begin{proof}
If \(V_\C\) is real or quaternionic, it must be
of the form \(V_\C = B \otimes U\), where \(U\) is an irreducible \(\C\)-local system.
However, since \(V_\C\) is isotrivial, \(U\) must have level zero. Since 
\(V_\C = V^{1,0} \oplus V^{0,1}\), \(B\) must have level one. In particular, it is 
at least two-dimensional.
We first claim that \(V_\C\) cannot be quaternionic. Indeed, if it were, then
by \Cref{proof-of-rep-classification}, \(h^0(\bigwedge^2 U)\) would contain
an alternating form. Consider the containment
\(\bigwedge^2 U \oplus \bigwedge^2 U \subseteq \bigwedge^2 V_\C\).
We know by Deligne's global invariant cycle theorem that 
\(h^0 \left(\bigwedge^2 V_\C\right) = 1\), so this containment is impossible, a contradiction.


Next, suppose that \(V_\C = B \otimes W_\C\) has real type. We claim
that the dimension of \(B\) is exactly two. Consider once again the containments
\begin{equation}
    \bigwedge^2 V_\C \supseteq (W_\C^{\otimes 2})^{\oplus \binom{\dim B}{2}}
    \supseteq (\text{Sym}^{2}W_\C)^{\oplus \binom{\dim B}{2}}
\end{equation}
Again by the global invariant cycle theorem, we know that the 
\(h^0(\bigwedge^2 V_\C) = 1\). But by \Cref{proof-of-rep-classification},
\(h^0(\text{Sym}^2 W_\C) = 1\). This forces \(\dim B = 2\).


Otherwise, suppose that \(V_\C\) is of complex type, 
\(V_\C = (B \otimes U) \oplus (\overline{B} \otimes \overline{U})\).
Then \(B\) and \(\overline{B}\) both have level zero. But now they must both be 
one-dimensional. If not, any subspace \(C \subseteq B\) is a Hodge subvariation,
with \((C \otimes U) \oplus (\overline{C} \otimes \overline{U}) \subsetneq V_\C\)
defined over \(\R\). This contradicts the irreducibility of \(V_\R\) as a Hodge 
structure. Therefore \(V_\C \simeq U \oplus \overline{U}\).


So far, we have not used the Hodge substructure \(H \subseteq V_\Q\). Consider 
the natural evaluation map 
\begin{equation}
    H \otimes_\Q \Hom_{\Q HS}(H, V_\Q) \to V_\Q.
\end{equation}
This map is always surjective. Suppose that \(E\) is not CM, i.e.
\(\End_{\Q HS}(H) = \Q\). Then by the previous theorem, 
\(\dim H \otimes_\Q \Hom_{\Q HS}(H, H^{\oplus n}) = 2n\), so that 
this map is an isomorphism of \(\Q\)-Hodge structures. Now notice that we may 
equip the domain with a Hodge representation, where \(\pi_1(B^\circ)\) acts 
trivially on \(H\) and \(\gamma . \varphi: h \to \varphi(\gamma h)\) for 
\(h \in H\) and \(\varphi \in \Hom_{\Q HS}(H, V_\Q)\). This extends the evaluation 
map to an isomorphism of Hodge representations. In particular, \(V_\C\) has real 
type.

Now suppose that \(E\) does have CM, \(\End_{\Q HS}(H) = K\) is an imaginary 
quadratic extension of \(\Q\). Since \(H_K = H_K^{1,0} \oplus H_K^{0,1}\) as a 
Hodge structure, the previous theorem gives us an analogous splitting of 
Hodge structures \(V_K = V_K^{1,0} \oplus V_K^{0,1}\). Moreover, Hodge automorphisms 
preserve the Hodge grading, so this gives a splitting of the Hodge representation 
\(V_K = V_K^{1,0} \oplus V_K^{0,1}\) into \(K\)-local systems. If they are 
non-isomorphic, then \(V_\C = V_\C^{1,0} \oplus V_\C^{0,1}\) is a splitting of 
\(V_\C\) into non-isomorphic \(\C\)-local systems. As we showed above, 
this is the complex type case, so that \(V_\C^{1,0}\) and \(V_\C^{0,1}\) are 
irreducible \(\C\)-local systems. 

Alternatively, suppose that \(V_K^{1,0} \simeq V_K^{0,1}\), so that 
\(V_K \simeq B \otimes U\) for \(U\) an irreducible \(K\)-local system.
If \(U = W_K\) for \(W\) a \(\Q\)-local system, then we are again in the real type
case. Thus \(B\) is two-dimensional and level one, so it is again identified 
with the Hodge substructure \(H \subseteq V_\Q\), and \(V_\Q = H \otimes W\).
Finally, suppose \(U\) is an irreducible \(K\)-local system which is not defined
over \(\Q\). Then \(U_\C\) is a \(\C\)-local system which is not defined 
over \(\R\). As we have seen, it cannot be quaternionic, so it must be of 
complex type. But if \(U_\C \neq \overline{U}_\C\), then \(U \neq \overline{U}\).
This contradicts our assumption.
\end{proof}

\end{document}
